\documentclass[11pt,a4paper]{article}
\usepackage[margin=2.4cm]{geometry}
\usepackage{amsmath,amssymb,amsthm}
\usepackage{booktabs}
\usepackage{array}
\usepackage{xcolor}
\usepackage[colorlinks=true,linkcolor=blue!50!black,urlcolor=blue!50!black]{hyperref}

\newcommand{\PB}{P^{\mathrm{B}}}
\newcommand{\PJ}{P^{\mathrm{J}}}
\newcommand{\assoc}{\operatorname{assoc}}
\newcommand{\OR}{\operatorname{OR}}

\theoremstyle{definition}
\newtheorem{definition}{Definition}
\theoremstyle{plain}
\newtheorem{proposition}{Proposition}
\theoremstyle{remark}
\newtheorem{remark}{Remark}

\newcommand{\PriorCells}{\frac{11}{60} & \frac{3}{20} \\ \frac{1}{15} & \frac{3}{5}}
\newcommand{\AbCells}{\frac{55}{98} & \frac{1}{14} \\ \frac{15}{98} & \frac{3}{14}}
\newcommand{\BaCells}{\frac{22}{61} & \frac{12}{305} \\ \frac{3}{11} & \frac{18}{55}}
\newcommand{\BmCells}{\frac{55}{94} & \frac{3}{47} \\ \frac{15}{94} & \frac{9}{47}}
\newcommand{\GeomRows}{prior $P$ & $\frac{1}{3}$ & $\frac{1}{4}$ & $\frac{1}{10}$ & $11$ \\
$\PJ_{AB}$ & $\frac{31}{49}$ & $\frac{5}{7}$ & $\frac{75}{686}$ & $11$ \\
$\PJ_{BA}$ & $\frac{2}{5}$ & $\frac{425}{671}$ & $\frac{72}{671}$ & $11$ \\
$\PB$ & $\frac{61}{94}$ & $\frac{35}{47}$ & $\frac{225}{2209}$ & $11$ \\}
\newcommand{\JAPzz}{\frac{11}{50}}
\newcommand{\JAPoz}{\frac{3}{50}}
\newcommand{\JAPcol}{\frac{7}{25}}
\newcommand{\ABzz}{\frac{55}{98}}
\newcommand{\Pzz}{\frac{11}{60}}
\newcommand{\Poz}{\frac{1}{15}}
\newcommand{\sAB}{\frac{3}{14}}
\newcommand{\sABn}{0.2143}
\newcommand{\sBA}{\frac{18}{55}}
\newcommand{\sBAn}{0.3273}
\newcommand{\sB}{\frac{9}{47}}
\newcommand{\sBn}{0.1915}
\newcommand{\uval}{- \frac{11}{188}}
\newcommand{\uvaln}{-0.0585}
\newcommand{\Kval}{\frac{468}{175}}
\newcommand{\Kvaln}{2.6743}
\newcommand{\DriftRows}{$\frac{1}{10}$ & $0.632653$ & $0.232653$ & $2.3265$ \\
$\frac{1}{100}$ & $0.426360$ & $0.026360$ & $2.6360$ \\
$\frac{1}{1000}$ & $0.402670$ & $0.002670$ & $2.6704$ \\
$\frac{1}{10000}$ & $0.400267$ & $0.000267$ & $2.6739$ \\}
\newcommand{\FHlo}{- \frac{1}{12}}
\newcommand{\FHhi}{\frac{1}{6}}
\newcommand{\FHfrac}{60}


\title{How a posterior is formed, and what ``pinning'' means\\
\large A worked example for \emph{Order-Sensitivity of Belief and Decision
Statistics under Jeffrey Conditioning}}
\author{}
\date{}

\begin{document}
\maketitle
\vspace{-2.2em}

\section{From prior to posterior}

The prior is a joint law on $\{0,1\}^{2}$, written in the coordinates
$(\alpha,\beta,c)$ --- the two marginals $P(A{=}0)=\alpha$, $P(B{=}0)=\beta$ and
the covariance $c$:
\begin{equation}
  P=\begin{pmatrix}\alpha\beta+c & \alpha(1-\beta)-c\\[2pt]
                   (1-\alpha)\beta-c & (1-\alpha)(1-\beta)+c\end{pmatrix},
  \qquad \assoc(P)=P_{00}P_{11}-P_{01}P_{10}=c .
  \label{eq:prior}
\end{equation}
Write $\alpha_{0}=\alpha,\ \alpha_{1}=1-\alpha$ and $\beta_{0}=\beta,\
\beta_{1}=1-\beta$.

A cue does not arrive as a likelihood. It arrives as a \emph{delivered
credence}: a target marginal $q=(q_{0},q_{1})$ on $A$'s partition, or
$r=(r_{0},r_{1})$ on $B$'s. A Jeffrey step installs that marginal and leaves
the conditionals alone:
\begin{equation}
  (J_{A}Q)(i,j)=q_{i}\,\frac{Q(i,j)}{Q(A{=}i)},
  \qquad
  (J_{B}Q)(i,j)=r_{j}\,\frac{Q(i,j)}{Q(B{=}j)} .
  \label{eq:step}
\end{equation}

\paragraph{The two reading orders.}
Composing two steps gives the two posteriors. Reading $A$ first and then $B$,
\begin{equation}
  \PJ_{AB}(i,j)
  \;=\;\bigl(J_{B}J_{A}P\bigr)(i,j)
  \;=\;r_{j}\,
    \frac{\dfrac{q_{i}P(i,j)}{\alpha_{i}}}
         {\displaystyle\sum_{k}\frac{q_{k}P(k,j)}{\alpha_{k}}},
  \label{eq:ab}
\end{equation}
and symmetrically, reading $B$ first,
\begin{equation}
  \PJ_{BA}(i,j)
  \;=\;\bigl(J_{A}J_{B}P\bigr)(i,j)
  \;=\;q_{i}\,
    \frac{\dfrac{r_{j}P(i,j)}{\beta_{j}}}
         {\displaystyle\sum_{l}\frac{r_{l}P(i,l)}{\beta_{l}}} .
  \label{eq:ba}
\end{equation}
\emph{The denominator is where the order lives.} The second step divides by the
marginal \emph{in force when it arrives}, and the first step has already moved
that marginal --- unless $c=0$, in which case the rows rescale proportionally
and nothing is disturbed.

\paragraph{The sequence-free benchmark.}
Re-express each delivered credence as a Bayes factor against the \emph{prior}
marginal, $\ell^{A}_{i}\propto q_{i}/\alpha_{i}$ and
$\ell^{B}_{j}\propto r_{j}/\beta_{j}$, and apply both at once:
\begin{equation}
  \PB(i,j)
  =\frac{P(i,j)\,\dfrac{q_{i}}{\alpha_{i}}\dfrac{r_{j}}{\beta_{j}}}
        {\displaystyle\sum_{k,l}P(k,l)\,\frac{q_{k}}{\alpha_{k}}\frac{r_{l}}{\beta_{l}}} .
  \label{eq:bench}
\end{equation}
A Bayes factor carries the same content whatever belief it meets, so
\eqref{eq:bench} does not care which cue came first. This is the reference
against which both routes are measured; it is not a rival updating rule, but
the same two impressions in an order-insensitive parametrisation.

\section{What ``pinning'' means}

\begin{definition}[pinned marginal]\label{def:pin}
A posterior's marginal on an attribute is \emph{pinned} by a cue if it equals
that cue's delivered credence \emph{exactly} --- as an identity in
$(\alpha,\beta,c,q_{0},r_{0})$, not merely to leading order in $c$.
\end{definition}

\begin{proposition}\label{prop:pin}
A Jeffrey step pins the marginal of the attribute it reads:
$(J_{A}Q)(A{=}i)=q_{i}$ for every $Q$. Hence the marginal of the cue read
\emph{last} is pinned, while the marginal of the cue read \emph{first} is not:
\[
  \PJ_{AB}(B{=}j)=r_{j},
  \qquad
  \PJ_{BA}(A{=}i)=q_{i},
  \qquad\text{identically in } c .
\]
\end{proposition}

\begin{proof}
Summing \eqref{eq:step} over $j$ gives
$\sum_{j}q_{i}Q(i,j)/Q(A{=}i)=q_{i}Q(A{=}i)/Q(A{=}i)=q_{i}$. The composite
claims follow because the final step of each route is the one that installs
that attribute's marginal.
\end{proof}

\begin{remark}[a plainer gloss: last writer wins]
Pinning is not an approximation or a coincidence --- it is the defining
property of the rule. What is at stake is which cue's credence \emph{survives}.
The second step rescales the other axis, and when $c\neq0$ the rows are not
proportional across columns, so it drags the first cue's marginal off the value
that cue delivered. The evaluator cannot honour both impressions at once:
whichever cue speaks last is taken at face value, and the earlier one has been
quietly overwritten through the covariance.
\end{remark}

\begin{remark}[what pinning is \emph{not}]
Pinning is a statement about \emph{marginals} only. It says nothing about the
association, which no attribute-local step can shift --- only rescale. The
next section makes both halves visible at once.
\end{remark}

\section{A worked example}

Take a prior with $\alpha=\tfrac13$, $\beta=\tfrac14$, $c=\tfrac1{10}$, and two
soft cues delivering $q=(\tfrac25,\tfrac35)$ on $A$ and $r=(\tfrac57,\tfrac27)$
on $B$. Substituting the three parameters into~\eqref{eq:prior} gives the prior;
feeding that prior and the two delivered credences through~\eqref{eq:ab},
\eqref{eq:ba} and~\eqref{eq:bench} gives the three posteriors:
\begin{alignat*}{2}
  P &= \begin{pmatrix}\PriorCells\end{pmatrix}
    &&\qquad\text{\small [prior parametrisation \eqref{eq:prior}]}\\[6pt]
  \PJ_{AB} &= \begin{pmatrix}\AbCells\end{pmatrix}
    &&\qquad\text{\small [$A$ read first, then $B$: \eqref{eq:ab}]}\\[6pt]
  \PJ_{BA} &= \begin{pmatrix}\BaCells\end{pmatrix}
    &&\qquad\text{\small [$B$ read first, then $A$: \eqref{eq:ba}]}\\[6pt]
  \PB &= \begin{pmatrix}\BmCells\end{pmatrix}
    &&\qquad\text{\small [Bayes-factor benchmark \eqref{eq:bench}]}
\end{alignat*}

\paragraph{One cell in full.}
To see~\eqref{eq:ab} at work, take the $(0,0)$ entry of the $A$-first route. The
inner Jeffrey step on $A$ rescales each row by $q_{i}/\alpha_{i}$, giving the
two entries of column $0$,
\[
  (J_{A}P)(0,0)=\frac{q_{0}P_{00}}{\alpha_{0}}
    =\frac{\tfrac25\cdot\Pzz}{\tfrac13}=\JAPzz,
  \qquad
  (J_{A}P)(1,0)=\frac{q_{1}P_{10}}{\alpha_{1}}
    =\frac{\tfrac35\cdot\Poz}{\tfrac23}=\JAPoz ,
\]
whose sum is the column mass $(J_{A}P)(B{=}0)=\JAPcol$ --- the denominator
of~\eqref{eq:ab}, and the quantity that already differs from $\beta=\tfrac14$
because the $A$-step has moved it. The outer step on $B$ then divides by that
column mass and multiplies by $r_{0}$:
\[
  \PJ_{AB}(0,0)
   = r_{0}\,\frac{(J_{A}P)(0,0)}{(J_{A}P)(B{=}0)}
   = \frac57\cdot\frac{\JAPzz}{\JAPcol}
   = \ABzz .
\]
Had $B$ been read first, that divisor would have been the \emph{prior} column
mass $\beta=\tfrac14$ instead, and the arithmetic would have landed elsewhere.
That single substitution is the whole of the order effect.

\paragraph{The geometry.}
All four beliefs live in the same simplex, so read each one off in the prior's
own coordinates --- two marginals and an association:

\begin{center}
\begin{tabular}{@{}lcccc@{}}
\toprule
belief & $P(A{=}0)$ & $P(B{=}0)$ & $\assoc$ & odds ratio \\
\midrule
\GeomRows
\bottomrule
\end{tabular}
\end{center}

\noindent Two things stand out.

\begin{enumerate}
\item \textbf{The pinned entries.} On route $AB$ the $B$-marginal is exactly
$r_{0}=\tfrac57$; on route $BA$ the $A$-marginal is exactly $q_{0}=\tfrac25$.
The \emph{other} marginal has been dragged off in each case, and the benchmark
pins neither.
\item \textbf{The odds ratio never moves.} It is $11$ in the prior and $11$ in
all three posteriors. Every route composes to a separable reweighting
$P(i,j)\mapsto g_{i}h_{j}P(i,j)/Z$, and such a reweighting multiplies
$\assoc$ by $g_{0}g_{1}h_{0}h_{1}/Z^{2}$ --- rescaling it, never shifting it.
The factor cancels in the ratio $P_{00}P_{11}/(P_{01}P_{10})$.
\end{enumerate}

\noindent So the order effect lives \emph{entirely in the marginals}. It cannot
build or erode the believed association at all.

\paragraph{The decision.}
Score the posterior against a threshold. Take a desirability vector $v$ paying
off only when both attributes are present, $s(Q)=\langle v,Q\rangle=Q_{11}$,
and a threshold $\tau=\tfrac14$:
\[
  s(\PJ_{AB})=\sAB\approx\sABn,\qquad
  s(\PJ_{BA})=\sBA\approx\sBAn,\qquad
  s(\PB)=\sB\approx\sBn .
\]
\begin{center}
\begin{tabular}{@{}lcc@{}}
\toprule
evaluator & score & action at $\tau=\tfrac14$ \\
\midrule
benchmark        & \sBn  & decline \\
reads $A$ first  & \sABn & decline \\
reads $B$ first  & \sBAn & \textbf{engage} \\
\bottomrule
\end{tabular}
\end{center}

\noindent Same prior, same two impressions, opposite hiring decision --- set
purely by reading order.

\paragraph{Who pays.}
The net surplus at the benchmark belief is
$u=s(\PB)-\tau=\uval\approx\uvaln$. The $B$-first reader flipped the action and
so bears the cost $|u|\approx\lvert\uvaln\rvert$. The $A$-first reader's
\emph{belief} also diverged from the benchmark --- but the action did not
change, so that evaluator costs nothing. Deviating in belief is not what is
expensive; flipping a decision is. This is why the loss concentrates in a thin
band at the threshold, and why it comes out an order smaller than the share of
evaluators affected.

\section{The drag is first order in \texorpdfstring{$c$}{c}}

The un-pinned marginal leaves its delivered credence at a rate given by
Proposition~C1$'$ of the manuscript,
\[
  K=\frac{q_{0}(1-q_{0})(r_{0}-\beta)}{\alpha\beta(1-\alpha)(1-\beta)}
   =\Kval\approx\Kvaln ,
\]
and the numerical drift converges to exactly that slope:

\begin{center}
\begin{tabular}{@{}cccc@{}}
\toprule
$c$ & $\PJ_{AB}(A{=}0)$ & drift from $q_{0}=\tfrac25$ & drift$/c$ \\
\midrule
\DriftRows
\bottomrule
\end{tabular}
\end{center}

\noindent At $c=0$ the drag vanishes entirely and both credences survive
intact: the two routes then coincide with each other and with the benchmark,
exactly.

\begin{remark}[a caution about the illustration]
The feasible range for $c$ at these marginals is
$[\FHlo,\FHhi]$ --- the Fr\'echet--Hoeffding bounds, forced by nonnegativity of
the four cells. The value $c=\tfrac1{10}$ used above is about \FHfrac\% of the
upper bound, and the first row of the table shows the consequence: a $0.23$
drift on a credence of $0.40$ is not a small correction. The figures are chosen
to make the mechanism visible, not to sit inside the regime where the
manuscript's expansions apply. Assumption~1 asks for $c$ small \emph{relative
to that bound}, and the lower rows of the table are where the paper's
first-order/second-order distinction becomes meaningful.
\end{remark}

\vfill
\noindent\rule{\linewidth}{0.4pt}\\
{\footnotesize Every number in this note was computed symbolically (exact
rationals, no floating point) with the verification suite in
\texttt{../sympy/}; the update formulas and Proposition~\ref{prop:pin} are
formalised in Lean in \texttt{../lean/}.}

\end{document}
